% Source: https://artofproblemsolving.com/community/c3h2522774_basic_trigonometry
% Stable identity: aops-topic-2522774-q01
\textbf{Estimate $\sqrt5+\sqrt[3]7$.} Since $2.2^2=4.84<5<5.0625=2.25^2$, we have $2.2<\sqrt5<2.25$. Since $1.9^3=6.859<7<7.077888=1.92^3$, we have $1.9<\sqrt[3]7<1.92$. So
\[4.1<\sqrt5+\sqrt[3]7<4.17.\]
Using $3.14<\pi<3.15$: $\pi<3.15<4.1$ and $\dfrac{3\pi}{2}>4.71>4.17$. Hence
\[\pi<\sqrt5+\sqrt[3]7<\dfrac{3\pi}2,\qquad\text{so}\qquad-\dfrac{3\pi}2<t<-\pi.\]

\textbf{Locate the point.} Turning by $t$ is the same as turning by $t+2\pi$, and $\dfrac\pi2<t+2\pi<\pi$. So $P_t$ lies in the $\boxed{\text{second quadrant}}$ ($x<0$, $y>0$).

\begin{asy}
usepackage("amsmath");
size(4cm);
defaultpen(fontsize(10pt));
draw((-1.35,0)--(1.35,0),gray,Arrow(5)); draw((0,-1.35)--(0,1.35),gray,Arrow(5));
draw(unitcircle,linewidth(0.8));
label("$x$",(1.35,0),E,fontsize(8pt)); label("$y$",(0,1.35),N,fontsize(8pt));
pair U(real t){return dir(t*180/pi);}
real t=-sqrt(5)-7^(1/3);
draw(arc((0,0),0.35,0,t*180/pi),red,Arrow(5));
dot(U(t)); draw((0,0)--U(t),dashed);
label("$P_t$",U(t),NW);
label("$t=-\sqrt5-\sqrt[3]7$",(0,-1.4),S,red+fontsize(8pt));
label("I",(0.75,0.75),fontsize(8pt)); label("II",(-0.75,0.75),fontsize(8pt)); label("III",(-0.75,-0.75),fontsize(8pt)); label("IV",(0.75,-0.75),fontsize(8pt));
\end{asy}
